\section{\S 3.4 沉淀溶解平衡}\label{sect:沉淀溶解平衡}

    \begin{itemize}
        \item \cemh{AgCl <=> Ag+ + Cl^-} 是哪种平衡？
        \item \cemh{NaCl_{(s)} <=> Na^+_{(aq)} + Cl^-_{(aq)}} 又是哪种平衡？
        \item \cemh{3Mg(OH)2 + 2FeCl3 \xleq{} 3MgCl2 + 2Fe(OH)3} 能发生吗？
        \item \cemh{Fe^{3+}} 在 \(\text{pH}=5\) 的溶液中能存在吗？
    \end{itemize}

    \textcolor{red}{我们如何判断某种金属正离子在 pH 一定的溶液中能否存在？}

    \subsection*{溶解平衡\textcolor{red}{——也是动态平衡}}\label{subsec:溶解平衡}

        \vspace{-\baselineskip}
            \[ \cemh{\text{固体溶质} <=>[\text{结晶}][\text{溶解}] \text{溶液中的溶质}}
            \quad \cemh{\text{气体分子} <=>[\text{析出}][\text{溶解}] \text{溶液中的微粒}} \]

        \begin{itemize}
            \item \(V_\text{溶解} > V_\text{结晶}\)\quad 溶质不断溶解减少（不饱和溶液）
            \item \(V_\text{溶解} = V_\text{结晶}\)\quad 溶质量不变，达到溶解平衡（饱和溶液）
            \item \(V_\text{溶解} < V_\text{结晶}\)\quad 溶质结晶析出，直至重新达到溶解平衡
        \end{itemize}

        溶解平衡的建立\textcolor{red}{——必须饱和溶液}

        \textcolor{blue}{\(\therefore\)难溶电解质的溶解平衡容易建立}

        难溶物：溶解度 \(< \SI{0.01}{g}/100\text{g 水}\)

        \textcolor{red}{难溶电解质，如何比较它们的溶解度大小？}

        \[ \cemh{AgCl_{(s)} <=>[\text{溶解}][\text{沉淀}] Ag^+_{(aq)} + Cl^-_{(aq)}} \]

        当\(v_\text{溶解} = v_\text{沉淀}\)时，实现平衡状态——沉淀溶解平衡

    \subsection[溶度积常数]{溶度积常数（简称：溶度积；符号：\(K_{\text{sp}}\)）}\label{subsec:溶度积常数}

        对于 \cemh{AgCl_{(s)}}，\( K_\text{sp} = \cemh{[Ag+]} \cdot \cemh{[Cl^-]} \quad K_\text{sp}(\cemh{AgCl}) = \cemh{[Ag+]} \cdot \cemh{[Cl^-]} \) 

        对于难溶电解质 \cemh{A_mB_n}，\cemh{A_mB_n_{(s)} <=> m A^{n+}_{(aq)} + n B^{m-}_{(aq)}}
        \[ K_\text{sp} = \cemh{[A^{n+}]}^m \cdot \cemh{[B^{m-}]}^n \]

        \subsubsection{溶度积与溶解度的关系}\label{subsubsec:溶度积与溶解度的关系}
            
            对于 \cemh{AB} 型难溶电解质，\( s = \sqrt{K_\text{sp}} \)\footnote{\(s_0(\si{\gram/100\gram\text{水}}) = \dfrac{s(\si{\mole\per\litre}) \cdot M}{10}\)}

            对于 \cemh{A2B} 型或 \cemh{AB2} 型难溶电解质，\( s = \sqrt[3]{\frac{K_\text{sp}}{4}} \)

            \textcolor{blue}{对于同类型难溶电解质，\(K_{\text{sp}}\)越小，在水中的溶解度也越小}
            
            \textcolor{blue}{不同类型，不能直接用\(K_{\text{sp}}\)比较，若相差很大，依然可以快速判断}

        \subsubsection{\texorpdfstring{\(K_{\text{sp}}\)}{Kₛₚ}的意义和影响因素}\label{subsubsec:Ksp意义和影响因素}
            
            \(K_{\text{sp}}\)的大小也反映了难溶电解质溶解趋势的大小，也反映了其生成沉淀的难易程度\textcolor{blue}{（越小越易生成）}
            
            \begin{xmp}
                [Ksp 意义]
                {}
                []
                []
                在含有 \cemh{Cl^-}、\cemh{Br-}、\cemh{I-} 的混合溶液中，滴加 \cemh{AgNO3} 溶液，最先生成 \cemh{AgI} 沉淀，最后生成 \cemh{AgCl} 沉淀。
            \end{xmp}

            \(K_{\text{sp}}\)的影响因素\(\left\{\begin{matrix} \makebox[1em][l]{难溶电解质本身的溶解度} \\
                \makebox[1em][l]{温度\textcolor{blue}{（大多数随温度升高而增大）}} \end{matrix}\right.\)

            \textcolor{red}{不受离子浓度影响}

        \subsubsection{溶度积规则}\label{subsubsec:溶度积规则}

            对于难溶电解质 \cemh{A_mB_n}，\cemh{A_mB_n_{(s)} <=> m A^{n+}_{(aq)} + n B^{m-}_{(aq)}}

            离子积\(Q\)，\(Q = c^m(\cemh{A^{n+}}) \cdot c^n(\cemh{B^{m-}})\)
            \begin{itemize}
                \item $Q < K_{sp}$时，无沉淀析出
                \item $Q = K_{sp}$时，沉淀溶解平衡状态
                \item $Q > K_{sp}$时，有沉淀析出，直至平衡状态
            \end{itemize}

        \begin{xmp}
            [沉淀生成判断]
            {}
            []
            []
            向 \SI{5}{\milli\litre} \SI{0.002}{\mole\per\litre} 的 \cemh{BaCl2} 溶液中加入 \SI{5}{\milli\litre} \SI{0.002}{\mole\per\litre} 的 \cemh{Na2SO4} 溶液，\Circled{1}计算判断是否有沉淀生成；\Circled{2}\cemh{Ba^{2+}}是否沉淀完全；\Circled{3}\cemh{Na2SO4} 溶液浓度改为\SI{0.0022}{\mole\per\litre}。\footnote{\textcolor{blue}{注：当溶液中离子浓度小于\textcolor{red}{\SI{E-5}{\mole\per\litre}} 时，即认为沉淀完全。}}
            
            （已知\(K_\text{sp}(\cemh{BaSO4}) = \num{1.08E-10}\)）
        \end{xmp}

        \begin{xmp}
            [同离子效应计算]
            {}
            []
            []
            计算 \cemh{BaSO4} 固体在下列条件下的溶解度（\si{\mole\per\litre}）\Circled{1}纯水中；\Circled{2}\SI{0.10}{\mole\per\litre} \cemh{Na2SO4} 溶液中；\Circled{3}\SI{0.001}{\mole\per\litre}的 \cemh{BaCl2} 溶液中。

            （已知\(K_\text{sp}(\cemh{BaSO4}) = \num{1.08E-10}\)）
            
            \Circled{1} \SI{1.04E-5}{\mole\per\litre}
            
            \Circled{2} \SI{1.08E-9}{\mole\per\litre}
            
            \Circled{3} \SI{1.08E-8}{\mole\per\litre}
            
        \end{xmp}

        \textcolor{blue}{同离子效应}：含有相同离子的易溶强电解质时，沉淀溶液平衡逆向移动，难溶电解质的溶解度减小。

        \textcolor{red}{\(\therefore\)沉淀剂往々会加过量}

        \subsubsection{难溶氢氧化物的开始沉淀 pH 和完全沉淀 pH}\label{subsubsec:难溶氢氧化物的开始沉淀pH和完全沉淀pH}

            金属离子的浓度为\SI{0.1}{\mole\per\litre}

            \begin{tabular}{|c|c|c|c|c|}
                \hline
                离子 & 氢氧化物 & 溶度积 \(K_{\text{sp}}\) & 开始沉淀 pH & 完全沉淀 pH \\
                \hline
                \cemh{Fe^{3+}} & \cemh{Fe(OH)3} & \num{2.79E-39} & 1.82& 2.82 \\
                \hline
                \cemh{Fe^{2+}} & \cemh{Fe(OH)2} & \num{4.87E-17} & 6.84 & 8.34 \\
                \hline
                \cemh{Cu^{2+}} & \cemh{Cu(OH)2} & \num{2.2E-20} & 5.17 & 6.67 \\
                \hline
                \cemh{Mg^{2+}} & \cemh{Mg(OH)2} & \num{8.45E-12} & 9.46 & 10.96 \\
                \hline
                \cemh{Al^{3+}} & \cemh{Al(OH)3} & \num{1.1E-33} & 3.68 & 4.68 \\
                \hline
            \end{tabular}

            \cemh{Fe(OH)3_{(s)} <=> Fe^{3+}_{(aq)} + 3OH^-_{(aq)}}

            \( \displaystyle K_\text{sp}(\cemh{Fe(OH)3}) = \cemh{[Fe^{3+}]} \cdot \cemh{[OH^-]}^3 \quad \cemh{[OH^-]} = \sqrt[3]{\frac{K_{\text{sp}}}{\cemh{[Fe^{3+}]}}} \quad \cemh{[H^+]} = \frac{K_\text{W}}{\cemh{[OH^-]}} \)

            \begin{xmp}
                [分离除杂应用]
                {}
                []
                []
                在酸化的 \cemh{CuSO4} 溶液中混有 \cemh{Fe^{3+}}，可加入\uline{\makebox[3em]{\cemh{CuO}}}来提高溶液的 pH，可以使 \cemh{Fe^{3+}}转变为 \cemh{Fe(OH)3} 而除去。该物质是否可以加过量？

                若还有 \cemh{Fe^{2+}}，可以先加入\uline{\makebox[3em]{\cemh{H2O2}}}，后再加入以上物质。

                \begin{enumerate}
                    \item \textcolor{red}{能与 \cemh{H+} 反应}
                    \item \textcolor{red}{不引入其他杂质离子\quad \(\therefore\)是含铜化合物}
                    \item \textcolor{red}{本身是难溶物，过量也可过滤除去}
                \end{enumerate}
                \begin{enumerate}
                    \item \textcolor{blue}{能够氧化 \cemh{Fe^{2+}}}
                    \item \textcolor{blue}{不引入杂质}
                \end{enumerate}
            \end{xmp}

    \subsection{沉淀溶解平衡的移动}\label{subsec:沉淀溶解平衡的移动}

        \subsubsection{常见金属离子沉淀物的颜色}\label{subsubsec:常见金属离子沉淀物的颜色}

            \begin{center}
            \small
            \begin{tabular}{|c|c|c|c|}
                \hline
                银盐或硫化物 & 颜色 & 氢氧化物 & 颜色 \\
                \hline
                氯化银 & 白色 & 氢氧化镁 & 白色 \\
                \hline
                溴化银 & 淡黄色 & 氢氧化铝 & 白色 \\
                \hline
                碘化银 & 黄色 & 氢氧化铜 & 黄色 \\
                \hline
                硫化银 & 黑色 & 氢氧化铁 & 红褐色 \\
                \hline
                硫化锌 & 白色 & 氢氧化亚铁 & 白色 \\
                \hline
            \end{tabular}
            \end{center}

        \subsubsection{沉淀的转化}\label{subsubsec:沉淀的转化}

            \makebox[20em][l]{\cemh{AgCl_{(s)} <=> Ag^+_{(aq)} + Cl^-_{(aq)}}} \(K_1 = K_\text{sp}(\cemh{AgCl})\)

            \makebox[20em][l]{\cemh{Ag^+_{(aq)} + I^-_{(aq)} <=> AgI_{(s)}}} \(K_2 = \frac{1}{K_\text{sp}(\cemh{AgI})}\)

            \cemh{AgCl_{(s)} + I^-_{(aq)} \xleq{} AgI_{(s)} + Cl^-_{(aq)}}

            \textcolor{red}{该过程的进行程度有多大？}

            \textcolor{blue}{方法一}：
            \[
                K = \frac{\cemh{[Cl^-]}}{\cemh{[I^-]}} = \frac{\cemh{[Cl^-]}\cdot\cemh{[Ag+]}}{\cemh{[I^-]}\cdot\cemh{[Ag+]}} = \frac{K_\text{sp}(\cemh{AgCl})}{K_\text{sp}(\cemh{AgI})} = \frac{\num{1.77E-10}}{\num{8.51E-17}} = \num{2.08E6}
            .\]

            \textcolor{blue}{方法二}：
            \[
                K = K_1 \times K_2 = \frac{K_\text{sp}(\cemh{AgCl})}{K_\text{sp}(\cemh{AgI})}
            .\]
            \textcolor{red}{\(\therefore\)该转化进行程度很大}


            \cemh{2AgI_{(s)} + S^{2-}_{(aq)} \xleq{} Ag2S_{(s)} + 2I^-_{(aq)}}

            \textcolor{red}{该过程的进行程度有多大？}

            \textcolor{blue}{方法一}：
            \[
                K = \frac{\cemh{[I^-]}^2}{\cemh{[S^{2-}]}} = \frac{\cemh{[I^-]}^2\cdot\cemh{[Ag+]}^2}{\cemh{[S^{2-}]}\cdot\cemh{[Ag+]}^2} = \frac{(K_\text{sp}(\cemh{AgI}))^2}{K_\text{sp}(\cemh{Ag2S})} = \frac{\left(\num{8.51E-17}\right)^2}{\num{6.69E-50}} = \num{1.08E17}
            .\]
            \textcolor{red}{\(\therefore\)该转化进行程度极大}

            \textcolor{blue}{方法二}：

            \makebox[20em][l]{\cemh{2AgI_{(s)} <=> 2Ag^+_{(aq)} + 2I^-_{(aq)}}} \(K_1 = (K_\text{sp}(\cemh{AgI})^2)\)

            \makebox[20em][l]{\cemh{2Ag^+_{(aq)} + S^{2-}_{(aq)} <=> Ag2S_{(s)}}} \(K_2 = \frac{1}{K_\text{sp}(\cemh{Ag2S})}\)

            \[
                K = K_1 \times K_2 = \frac{(K_\text{sp}(\cemh{AgI}))^2}{K_\text{sp}(\cemh{Ag2S})}
            .\]

            \textcolor{blue}{一般来说，同种离子的难溶物，溶解度小的易转化成溶解度更小的；\(K_\text{sp}\)小相差越大，转化越完全。}


            \cemh{3Mg(OH)2 + 2FeCl3 \xleq{} 3MgCl2 + 2Fe(OH)3}

            \textcolor{red}{这个反应能发生吗？进行程度大吗？}

            \[
                K = \frac{(K_\text{sp}(\cemh{Mg(OH)2}))^3}{(K_\text{sp}(\cemh{Fe(OH)3}))^2} = \frac{\left(\num{8.45E-12}\right)^3}{\left(\num{2.79E-39}\right)^2} = \num{7.75E43}
            .\]

            \textcolor{red}{【归纳】}

            沉淀转化反应的\(K = \frac{\left(\text{前一个沉淀的}K_\text{sp}\right)^\text{幂次方}}{\left(\text{后一个沉淀的}K_\text{sp}\right)^\text{幂次方}}\)
            

            \begin{xmp}
                [硫酸亚铁混有亚铜离子]
                {应用实例}
                []
                []
                硫酸亚铁溶液中混有 \cemh{Cu2+}，如何除去？

                可以通过加入 \cemh{FeS} 充分反应后过滤除去。

                \[
                    \cemh{FeS_{(s)} + Cu^{2+}_{(aq)} \xleq{} CuS_{(s)} + Fe^{2+}_{(aq)}}
                ,\]

                \[
                    K = \frac{K_\text{sp}(\cemh{FeS})}{K_\text{sp}(\cemh{CuS})} = \frac{\num{1.59E19}}{\num{1.27E36}} = \num{1.25E17}
                .\]
            \end{xmp}

            \textcolor{blue}{沉淀转化当\(K_\text{sp}\)接近时，则取决于离子浓度}

            \begin{xmp}
                [Ksp 接近时的判据]
                {}
                []
                []
                \cemh{BaSO4_{(s)} + CO3^{2-}_{(aq)} <=> BaCO3_{(s)} + SO4^{2-}_{(aq)}}

                \[
                    K = \frac{\cemh{[SO4^{2-}]}}{\cemh{[CO3^{2-}]}} = \frac{K_\text{sp}(\cemh{BaSO4})}{K_\text{sp}(\cemh{BaCO3})} = \frac{\num{1.08E-10}}{\num{8.1E-9}} = \frac{1}{75}
                .\]

                \(\therefore\)当\(c(\cemh{CO_3^{2-}}) > 75c(\cemh{SO_4^{2-}})\)时，\(Q = \frac{c{\cemh{SO_4^{2-}}}}{c{\cemh{CO_3^{2-}}}} < \frac{1}{75} = K\)，反应正向进行。
            \end{xmp}

        \subsubsection{沉淀的溶解}\label{subsubsec:沉淀的溶解}

            \noindent\textbf{（1）氢氧化物溶于酸（生成水）}
            
            \cemh{Mg(OH)2_{(s)} <=> Mg^{2+}_{(aq)} + 2OH^-_{(aq)}}

            \cemh{2H^+_{(aq)} + 2OH^-_{(aq)} \xleq{} 2H2O_{(l)}}
            
            \cemh{Mg(OH)2_{(s)} + 2H^+_{(aq)} \xleq{} Mg^{2+}_{(aq)} + 2H2O_{(l)}}
            
            \textcolor{red}{计算常温下该反应的\(K\)}
            \[
                K = \frac{K_\text{sp}(\cemh{Mg(OH)2})}{K_\text{W}^2} = \frac{8.45E-12}{\left(1.0E-14\right)^2} = \num{8.45E16}
            .\]

            \noindent\textbf{（2）碳酸盐溶于酸（生成气体）}

            \cemh{CaCO3_{(s)} + 2H^+_{(aq)} <=> Ca^{2+}_{(aq)} + H2CO3_{(aq)}}

            \textcolor{red}{计算常温下该反应的\(K\)}
            \[
                K = \frac{K_\text{sp}(\cemh{CaCO3})}{K_\text{a1}(\cemh{H2CO3})\cdot K_\text{a2}(\cemh{H2CO3})} = \frac{4.96E-9}{\num{4.2E-7}\cdot\num{4.8E-11}} = 2.46E8
            .\]
            \cemh{H2CO3_{(aq)} \xleq{} H2O_{(l)} + CO2_{(g)}}
            
            再考虑 \cemh{CO2_{(g)}}的逸出，实际\(K\)还要大。

            \noindent\textbf{（3）氢氧化镁溶于浓氯化铵（生成弱电解质）}

            \cemh{Mg(OH)2_{(s)} + 2NH4^+_{(aq)} <=> Mg^{2+}_{(aq)} + 2NH3 . H2O_{(aq)}}

            \[
                K = \frac{K_\text{sp}(\cemh{Mg(OH)2})}{(K_\text{b}(\cemh{NH3 . H2O}))^2} = \frac{\num{8.45E-12}}{(\num{1.8E-5})^2} = 0.026
            .\]

            \textcolor{red}{\(\therefore\)需要浓的氯化铵溶液才能溶解氢氧化镁}

            \noindent\textbf{（4）氯化银溶于氨水（生成配合物）}

            \cemh{AgCl_{(s)} + 2NH3_{(aq)} <=> [Ag(NH3)2]^+_{(aq)} + Cl^-_{(aq)}}

            \makebox[20em][l]{\cemh{AgCl_{(s)} <=> Ag^+_{(aq)} + Cl^-_{(aq)}}} \(K_\text{sp}(\cemh{AgCl})\)

            \makebox[20em][l]{\cemh{Ag^+_{(aq)} + 2NH3_{(aq)} <=> [Ag(NH3)2]^+_{(aq)}}} \(K_\text{稳}(\cemh{[Ag(NH3)2]+})\)

            \noindent\textbf{（5）难溶金属硫化物是否都溶于强酸}

            \cemh{FeS_{(s)} + 2H^+_{(aq)} <=> H2S_{(aq)} + Fe^{2+}_{(aq)}}

            \[
                K = \frac{K_\text{sp}(\cemh{FeS})}{K_\text{a1}(\cemh{H2S})\cdot K_\text{a2}(\cemh{H2S})} = \frac{\num{1.59E-19}}{\num{5.7E-8}\cdot\num{1.2E-15}} = \num{2.32E3}
            .\]

            \cemh{CuS_{(s)} + 2H^+_{(aq)} <=> H2S_{(aq)} + Cu^{2+}_{(aq)}}

            \[
                K = \frac{K_\text{sp}(\cemh{CuS})}{K_\text{a1}(\cemh{H2S})\cdot K_\text{a2}(\cemh{H2S})} = \frac{\num{1.27E-36}}{\num{5.7E-8}\cdot\num{1.2E-15}} = \num{1.86E-14}
            .\]

            \textcolor{red}{\(\therefore\)\cemh{FeS} 溶于强酸，\cemh{CuS} 不溶}

    \subsection*{如何衡量判断竞争反应}\label{subsec:竞争反应判断}
    \addcontentsline{toc}{subsection}{如何衡量判断竞争反应}
        \textcolor{red}{——利用常数}
